\subsection{Camada de dados e rastreabilidade}\label{sec:dados_rastreabilidade}
O DoseCerta utiliza dois níveis de persistência.  Na borda, a central guarda em SQLite o que precisa para operar sem rede: configuração ativa, agenda, decisões de dose, execuções e as filas de entrada e saída. No servidor, o PostgreSQL concentra o histórico compartilhado, os cadastros e as regras de acesso. Os dois níveis se comunicam apenas pela API; o aparelho não acessa o banco central.

A escolha de manter dois bancos decorre do requisito de operação local (RNF01) e do requisito de não repetir movimento (RNF02). O banco local é projetado para conservar o registro de uma operação em andamento diante de quedas de energia; a eficácia dessa recuperação depende de ensaio controlado; o banco central precisa reunir os eventos de várias fontes e preservar sua origem para consulta no aplicativo.

\subsubsection{Banco local da central}
O banco local foi configurado para operar em modo WAL (\emph{write-ahead logging}) com sincronização completa a cada transação (\texttt{synchronous = FULL}), chaves estrangeiras ativas e verificação de integridade na inicialização. A Tabela~\ref{tab:dados-local} resume suas tabelas. As restrições de domínio ficam no próprio esquema: posição do módulo entre 1 e 2, encaixe entre 1 e 4, quantidade por dose positiva, intervalo mínimo não negativo e janelas de tolerância coerentes.

{\small\setlength{\tabcolsep}{3pt}\begin{longtable}{p{\dimexpr0.24\linewidth-2\tabcolsep\relax}p{\dimexpr0.66\linewidth-2\tabcolsep\relax}}
\caption{Tabelas do banco local da central (SQLite)}\label{tab:dados-local}\\
\hline\textbf{Tabela} & \textbf{Função}\\\hline\endfirsthead
\hline\textbf{Tabela} & \textbf{Função}\\\hline\endhead
\texttt{config\_snapshots} & Versões da configuração recebida do servidor; somente uma ativa por vez.\\\hline
\texttt{schedules} & Agenda local: horários fixos ou sob demanda, posição, dose, intervalo mínimo e tolerâncias.\\\hline
\texttt{dose\_decisions} & Cada decisão de liberar ou bloquear uma dose, com chave única por pedido e motivo.\\\hline
\texttt{dose\_executions} & Execução física de uma decisão: iniciada, concluída, falha ou \emph{incerta}.\\\hline
\texttt{inbox} & Comandos recebidos do servidor, com estado de processamento e expiração.\\\hline
\texttt{outbox} & Eventos a enviar ao servidor, com chave de idempotência única e contagem de tentativas.\\\hline
\texttt{runtime\_state} & Estado persistente do serviço entre reinícios.\\\hline
\end{longtable}}{\small\noindent Fonte: Autoria própria (2026).\par}

O estado \emph{incerto} de uma execução é o ponto central da recuperação. Se a energia cai ou a comunicação com o submódulo se perde durante um giro, a execução não é marcada como concluída nem como falha. Ela permanece incerta, e a central não aceita novo movimento naquela posição até que a situação seja reconciliada com a leitura do submódulo ou com a confirmação humana. Esse comportamento é esperado na política rigorosa; o ambiente de apresentação admite exceções e deve ser avaliado separadamente. Registros de recusa e de reconciliação não provam a recuperação física em toda falha de energia.

\subsubsection{Banco central}
O esquema organiza as entidades por função e utiliza vínculos, restrições, políticas de acesso por linha e gatilhos. A arquitetura desses controles não constitui auditoria de segurança; privilégios da API podem contornar RLS. A Tabela~\ref{tab:dados-central} agrupa as tabelas por função.

{\small\setlength{\tabcolsep}{3pt}\begin{longtable}{p{\dimexpr0.2\linewidth-2\tabcolsep\relax}p{\dimexpr0.4\linewidth-2\tabcolsep\relax}p{\dimexpr0.3\linewidth-2\tabcolsep\relax}}
\caption{Tabelas do banco central agrupadas por função}\label{tab:dados-central}\\
\hline\textbf{Grupo} & \textbf{Tabelas} & \textbf{Função}\\\hline\endfirsthead
\hline\textbf{Grupo} & \textbf{Tabelas} & \textbf{Função}\\\hline\endhead
Pessoas e acesso & \texttt{patient}, \path{patient_responsible}, \texttt{app\_user\_profile}, \path{device_caregiver_access} & Pacientes, responsáveis e quem pode ver o quê\\\hline
Tratamento & \texttt{medication}, \path{medication_schedule}, \texttt{intake\_event} & Medicamentos, horários e registros de tomada\\\hline
Rotina de cuidado & \texttt{care\_activity}, \texttt{care\_observation}, \texttt{care\_review}, \texttt{care\_preference}, \texttt{care\_notice\_read} & Atividades, relatos e leitura de avisos\\\hline
Aparelho & \texttt{device}, \texttt{device\_level}, \texttt{device\_module}, \path{device_compartment}, \path{device_compartment_inventory}, \path{device_patient_assignment} & Topologia física, posições, estoque e vínculo com pacientes\\\hline
Identidade do aparelho & \texttt{device\_credential}, \path{device_claim_token}, \path{device_config_revision} & Credenciais, pareamento e versões de configuração\\\hline
Operação e auditoria & \texttt{device\_command}, \texttt{device\_status}, \texttt{domain\_event}, \texttt{api\_idempotency} & Comandos, sinais de vida, eventos e chaves de idempotência\\\hline
Simulação & \path{device_simulation_profile}, \path{device_simulation_action}, \path{device_simulation_artifact} & Aparelho simulado para desenvolvimento\\\hline
\end{longtable}}{\small\noindent Fonte: Autoria própria (2026).\par}

Três tabelas sustentam a rastreabilidade. A \texttt{device\_command} guarda cada comando enviado à central, com quem o pediu, quando foi entregue, reconhecido e aplicado, e o resultado devolvido pelo aparelho. A \texttt{domain\_event} registra cada fato relevante (comando enfileirado, dose dispensada, saída confirmada, configuração publicada) com o tipo de ator que o produziu: o aparelho, um operador humano ou o próprio sistema. A \texttt{api\_idempotency} associa cada escrita a uma chave única, para reconhecer um reenvio da mesma operação. A garantia depende de preservar a chave e a semântica da operação; não se presume para requisições com chaves diferentes ou para efeitos físicos sem confirmação.

\subsubsection{Dois ambientes}
O servidor mantém dois ambientes com bancos separados no mesmo PostgreSQL. O ambiente \emph{oficial} aplica a política de travas \emph{rigorosa}: o aparelho só executa um movimento quando todas as pré-condições são satisfeitas. O ambiente \emph{simulado} aplica uma política de \emph{apresentação}, com travas mínimas, e só vale quando o aparelho e o servidor concordam com ela. A troca de ambiente é feita na central por um procedimento administrativo que altera o endereço da API e a política, recusa a troca com o motor em movimento e não toca na credencial, no banco local nem nas redes Wi-Fi.

Essa separação permitiu apresentar o protótipo sem misturar dados de demonstração com os registros oficiais. Em contrapartida, os resultados precisam indicar em qual ambiente foram obtidos, pois a política de apresentação pode fechar um ensaio automaticamente, sem a contagem humana exigida na avaliação.

\subsubsection{Privacidade e alcance dos controles}
A implementação utiliza HTTPS, tokens de sessão, credenciais próprias para a central, pareamento com token de uso único armazenado como resumo criptográfico, desafio de autenticação no Bluetooth e verificação de integridade e assinatura das releases da Raspberry. A descrição dos controles não equivale a uma auditoria independente. Para as ESP32 permanece a distinção da Seção~\ref{cap:embarcados}: SHA-256 não é assinatura criptográfica. 

A existência de políticas RLS, chaves estrangeiras e restrições não garante sozinha autorização por pessoa. A configuração apresentada é limitada a um paciente por cadastro; avaliação multiusuário e multipaciente é uma possibilidade de continuidade, não critério de aceite desta versão. A configuração acadêmica utiliza conta operadora compartilhada, privilégios da API que podem contornar RLS e credencial estática do MCP. As três visões da interface não foram interpretadas como três identidades autenticadas independentes. 

O PostgreSQL no servidor guarda cadastros, agenda, relatos, comandos, eventos e telemetria, com bancos separados para Oficial e Simulado. O SQLite da central guarda configuração, agenda e filas locais; os resultados detalhados da visão permanecem na central, enquanto ao servidor chegam pedido e estado. Conversas e imagens de receitas enviadas ao assistente são processadas no Azure/Microsoft Foundry. Portanto, processamento local da visão não significa processamento exclusivamente local de todas as imagens do sistema. 

Retenção, exclusão, restauração, acesso a observações reservadas e tratamento de dados no serviço externo dependem de política e comprovação específicas. Não se afirma conformidade legal, anonimização integral ou proteção contra toda repetição de movimento. A idempotência dos registros e a contagem física do coletor permanecem avaliações diferentes.
